\documentclass[a4paper, 12pt, landscape]{scrartcl}

\usepackage[a4paper, landscape, margin=1.0cm, top=1.2cm, bottom=1.0cm]{geometry}
\usepackage[ngerman]{babel}
\usepackage{fontspec}
\usepackage{xcolor}
\usepackage[most]{tcolorbox}

\setmainfont{Noto Sans}
\pagestyle{empty}

% Silbentrennung komplett ausschalten für sauberen Plakatdruck
\hyphenpenalty=10000
\exhyphenpenalty=10000

\definecolor{primaryBlue}{RGB}{20, 50, 100}
\definecolor{accentOrange}{RGB}{220, 100, 30}
\definecolor{textDark}{RGB}{15, 23, 42}
\definecolor{subtleGray}{RGB}{100, 116, 139}

\newcommand{\aussagePlakat}[5]{%
    \noindent\begin{tcolorbox}[
        enhanced,
        colback=white,
        colframe=primaryBlue,
        boxrule=3pt,
        arc=4mm,
        width=\textwidth,
        height=17.2cm,
        valign=center,
        halign=center,
        top=0.6cm, bottom=0.6cm, left=1.0cm, right=1.0cm,
        title={\LARGE\bfseries\color{white}\quad #1 \quad},
        attach boxed title to top center={yshift=-4.5mm},
        boxed title style={colback=primaryBlue, arc=2mm, boxrule=0pt, left=10mm, right=10mm, top=2.8mm, bottom=2.8mm}
    ]
        \vspace{-0.1cm}
        {\fontsize{#2}{#3}\selectfont\bfseries\color{textDark}
        \glqq #4\grqq}
        
        \vfill
        
        \begin{tcolorbox}[
            enhanced,
            colback=white,
            colframe=primaryBlue!30,
            boxrule=1pt,
            arc=2mm,
            width=0.96\textwidth,
            halign=center,
            top=2.5mm, bottom=2.5mm
        ]
            {\normalsize\bfseries\color{primaryBlue}Philosophie Q1 \enskip$\cdot$\enskip Herr Dayi} 
            \quad\textbar\quad 
            {\normalsize\bfseries\color{subtleGray}#5}
        \end{tcolorbox}
    \end{tcolorbox}
    \clearpage
}

\begin{document}

% === AUSSAGE A ===
\aussagePlakat{AUSSAGE A}{26pt}{38pt}{%
Das Tier produziert nur unter der Herrschaft des physischen Bedürfnisses,\\
während der Mensch frei produziert\\
und erst wahrhaft produziert in der Freiheit [\dots]\\
der Mensch formiert daher auch nach den Gesetzen der Schönheit.%
}{Karl Marx: Ökonomisch-philosophische Manuskripte (1844)}

% === AUSSAGE B ===
\aussagePlakat{AUSSAGE B}{28pt}{42pt}{%
Was das Produkt seiner Arbeit ist, ist er nicht.\\
Je mehr Gegenstände der Arbeiter produziert,\\
um so mehr gerät er unter die Herrschaft seines Produkts,\\
des Kapitals.%
}{Karl Marx: Ökonomisch-philosophische Manuskripte (1844)}

% === AUSSAGE C ===
\aussagePlakat{AUSSAGE C}{26pt}{38pt}{%
Seine Arbeit ist daher nicht freiwillig,\\
sondern gezwungen, Zwangsarbeit.\\
Sie ist nicht die Befriedigung eines Bedürfnisses,\\
sondern nur ein Mittel,\\
um Bedürfnisse außer ihr zu befriedigen.%
}{Karl Marx: Ökonomisch-philosophische Manuskripte (1844)}

% === AUSSAGE D ===
\aussagePlakat{AUSSAGE D}{27pt}{40pt}{%
Es kommt daher zum Resultat,\\
dass der Mensch sich nur mehr in seinen tierischen Funktionen\\
(Essen, Trinken, Zeugen) als freitätig fühlt\\
und in seinen menschlichen Funktionen nur mehr als Tier.%
}{Karl Marx: Ökonomisch-philosophische Manuskripte (1844)}

\end{document}
